% !TeX root = ../教材题库_全部条目_转写工作稿.tex
% Source: XJ-B1-C06-SU-P263-Y0217
% 第6章 / 小结与复习 小结与复习 / 教材 p.263 / PDF 第270页
% OCR未校对：公式与文字均可能有误；未完成数学TeX转写；禁止直接发给学生
% 原题及相邻原解截图（核对用）：../题目原图/XJ-B1-C06-SU-P263-Y0217-01.webp
% 本文件只在转写完成后作为题干使用；OCR 条目尚不是可靠题干。
\begin{bankitem}{第6章 小结与复习 第4题}{第6章 / 小结与复习 小结与复习 / 教材 p.263 / PDF 第270页}{XJ-B1-C06-SU-P263-Y0217}{OCR未校对：公式与文字均可能有误；未完成数学TeX转写；禁止直接发给学生}
\begin{ocrdraft}
[教材p.263 / PDF第270页]
\par
4．用样本估计总体是统计的基本思想．具体而言，本章分别从以下几个维度
\par
来介绍：
\par
（D用样本的数字特征估计总体的数字特征．你能说出平均数、中位数和众
\par
数，以及极差、方差与标准差的统计含义吗？它们在统计中起到的作用是什么？
\par
（2）用样本的频率分布直方图来刻画一组数据的分布情况，进而估计总体分
\par
布，这将有助于我们从整体上把握研究对象．
\par
（3）用样本估计百分位数．百分位数用于描述总体或样本的百分位点，多个
\par
百分位数结合使用可更全面地描述数据的分布特征．
\end{ocrdraft}
\end{bankitem}
